\documentclass[12pt,a4paper]{article}

% ── Packages ─────────────────────────────────────────────────────────────────
\usepackage[margin=2.5cm]{geometry}
\usepackage{amsmath,amssymb,amsthm}
\usepackage{graphicx}
\usepackage{booktabs}
\usepackage{hyperref}
\usepackage{xcolor}
\usepackage{caption}
\usepackage{subcaption}
\usepackage{float}
\usepackage{microtype}
\usepackage[numbers,sort&compress]{natbib}

\hypersetup{
  colorlinks=true, linkcolor=blue!60!black,
  citecolor=green!50!black, urlcolor=blue!70!black,
}

\newcommand{\norm}[1]{\left\lVert#1\right\rVert}
\newcommand{\Lap}{\Delta}
\newcommand{\lapinv}{\Delta^{-1}}
\newcommand{\Real}{\mathbb{R}}
\newcommand{\Torus}{\mathbb{T}^2}
\newcommand{\Proj}{P}
\newcommand{\Hmone}{H^{-1}}
\newcommand{\half}{\tfrac{1}{2}}
\newcommand{\dx}{\,dx\,dy}
\newcommand{\Dop}[1]{D_{#1}}
\newcommand{\Linv}{\Lambda^{-1}}

\theoremstyle{definition}
\newtheorem{remark}{Remark}
\theoremstyle{plain}
\newtheorem{proposition}{Proposition}

% ── Title ─────────────────────────────────────────────────────────────────────
\title{%
  \textbf{Optimal Mixing of Passive Scalars}\\[6pt]
  \large A Pseudo-Spectral Numerical Study of the\\
  Lin--Thiffeault--Doering Velocity
}
\author{Adebanji Adelowo}
\date{April 27, 2021}

% ─────────────────────────────────────────────────────────────────────────────
\begin{document}
\maketitle

\begin{abstract}
\noindent
This thesis investigates the optimal mixing of passive scalars in
incompressible fluid flow using the instantaneous-optimal (greedy) velocity
field derived by Lin, Thiffeault \& Doering (2011). The central objective is
to numerically characterise the decay rate of the $H^{-1}$ mix norm under
the LTD optimal stirring strategy and to examine how the mixing timescale
depends on the spatial scale of the initial scalar distribution.

The scalar transport equation is solved numerically on the two-dimensional
torus $[0,1]^2$ using a pseudo-spectral discretisation (fast Fourier
transforms in space) together with an adaptive Dormand--Prince RK45 scheme
for time integration. Four families of initial conditions -- parameterised
by a support-size parameter $a$ -- are studied, with $L^p$ norm conservation
used as a resolution diagnostic and stopping criterion.

Numerical results confirm that, under the dynamically chosen LTD velocity,
the $H^{-1}$ mix norm decays approximately exponentially, consistent with
the robust exponential decay reported numerically by Lin, Thiffeault and
Doering for the same steepest-descent strategy. Writing this decay as
$\norm{\theta(t)}_{H^{-1}} \approx \norm{\theta_0}_{H^{-1}} e^{-rt}$ with
decay-rate magnitude $r>0$, a representative case ($a = 0.5$) gives
$r \approx 0.44$ (mixing timescale $\tau = 1/r \approx 2.27$). Across eight
values of $a \in [0.5,\, 0.9375]$, the fitted decay-rate magnitude scales
approximately as $r \propto a^{-1.78}$, appreciably steeper than the
$a^{-1}$ dependence characteristic of the Iyer--Kiselev--Xu
enstrophy-constrained lower bound. Comparing two resolutions ($N=32$:
exponent $\approx -1.61$; $N=64$: exponent $\approx -1.78$) shows this
fitted exponent is resolution-dependent and not yet numerically converged;
since the higher-resolution estimate moves farther from, not closer to,
the benchmark $-1$ scaling, these two resolutions alone cannot establish
whether the discrepancy reflects a finite-resolution artefact, a genuine
feature of the greedy strategy, or both.

These results provide a foundation for future studies at higher spectral
resolution and for comparison with non-greedy global optimisation strategies.

\medskip\noindent\textbf{Keywords:} passive scalar mixing, $H^{-1}$ mix norm,
pseudo-spectral methods, optimal stirring, Lin--Thiffeault--Doering velocity,
incompressible fluid mechanics
\end{abstract}

\tableofcontents
\bigskip

% ─────────────────────────────────────────────────────────────────────────────
\section{Introduction}
\label{sec:intro}

This report documents a pseudo-spectral numerical simulation of the
\emph{optimal mixing problem} for passive scalars on the two-dimensional
torus $\mathbb{T}^2 = [0,1]^2$.
The simulation reproduces and extends Gautam Iyer's earlier MATLAB
implementation of the instantaneous-optimal (steepest-descent) stirring
strategy derived by Lin, Thiffeault and Doering \cite{LTD2011}.

\subsection{Motivation}

Mixing a substance efficiently under a fixed stirring budget is a question
that appears across scales: from folding cream into coffee, to homogenising
reactants in a chemical reactor, to the turbulent dispersal of pollutants
and phytoplankton in the ocean. In all of these settings the underlying
physics is the same passive-scalar transport equation studied here, and the
same question recurs -- given a fixed amount of ``stirring energy,'' what
motion mixes fastest? The LTD framework answers this question precisely for
an idealised, instantaneous notion of optimality, and is a natural entry
point into the broader mathematical theory of mixing, which connects
functional analysis (Sobolev norms as mixing diagnostics), the calculus of
variations (constrained optimisation of the velocity field), and dynamical
systems (chaotic advection, filamentation, and the eventual limits imposed
by finite resolution or molecular diffusion).

\subsection{Goals}
\begin{itemize}
  \item Reproduce the results of Gautam Iyer's earlier MATLAB implementation
        of this simulation.
  \item Demonstrate that the $H^{-1}$ mix norm decays exponentially under the
        Lin--Thiffeault--Doering (LTD) optimal velocity.
  \item Characterise how the mixing timescale depends on the support
        size parameter~$a$.
  \item Derive the LTD velocity rigorously, rather than simply quoting it,
        so that the numerical implementation can be checked term-by-term
        against the theory (Section~\ref{sec:math}).
  \item Compare the fitted decay-rate scaling in $a$ against the benchmark
        $a^{-1}$ dependence of the Iyer--Kiselev--Xu enstrophy-constrained
        lower bound, and test whether the fitted exponent is
        resolution-converged via a direct two-resolution comparison
        (Section~\ref{sec:resolution}).
\end{itemize}

\subsection{Structure of this Report}

Section~\ref{sec:math} develops the mathematical framework, including a
full derivation of the LTD optimal velocity via constrained optimisation.
Section~\ref{sec:numerics} documents the pseudo-spectral numerical scheme.
Section~\ref{sec:idata} describes the initial data families.
Section~\ref{sec:results} presents the main simulation results at the
published resolution $N=64$.
Section~\ref{sec:resolution} adds a resolution-sensitivity study not present
in Gautam Iyer's earlier MATLAB implementation.
Section~\ref{sec:conc} closes with conclusions.

% ─────────────────────────────────────────────────────────────────────────────
\section{Mathematical Framework}
\label{sec:math}

\subsection{The Passive Scalar Equation}

We study the advection of a passive scalar $\theta : \Torus\times[0,\infty)
\to \Real$ governed by

\begin{equation}
  \label{eq:pde}
  \partial_t\theta + (u\cdot\nabla)\theta = 0,
  \qquad \theta(\cdot,0) = \theta_0,
\end{equation}

where $u:\Torus\to\Real^2$ is a prescribed divergence-free velocity field,
$\nabla\cdot u = 0$. The scalar $\theta$ is \emph{passive}: it is transported
by $u$ but exerts no back-reaction on the flow, and there is no molecular
diffusion. Equation~\eqref{eq:pde} states that $\theta$ is constant along
particle trajectories $\dot X(t) = u(X(t),t)$; since the flow map generated
by a divergence-free $u$ is measure-preserving, every level set $\{\theta =
c\}$ is merely rearranged in time, never created or destroyed. This is the
single physical fact underlying everything that follows: mixing is a purely
kinematic phenomenon of rearrangement, and the question of ``optimal
mixing'' is really a question about the most effective way to rearrange.

\subsection{The Mix Norm}
\label{sec:mixnorm-def}

The quality of mixing is quantified by the \emph{$H^{-1}$ mix norm}:

\begin{equation}
  \label{eq:mixnorm}
  \norm{\theta}_{H^{-1}}^2
  = \sum_{k\neq 0} \frac{|\hat\theta(k)|^2}{4\pi^2|k|^2},
\end{equation}

where $\hat\theta(k)$ denotes the $k$-th Fourier coefficient of $\theta$.
Equivalently, on the mean-zero subspace, $\norm{\theta}_{H^{-1}}^2 =
-\int_\Torus \theta\,\lapinv\theta \dx$: writing $\phi = \lapinv\theta$
(i.e.\ $\Lap\phi=\theta$), integration by parts gives $-\int\theta\phi\dx =
-\int(\Lap\phi)\phi\dx = \int|\nabla\phi|^2\dx \ge 0$, confirming that
\eqref{eq:mixnorm} really is a squared norm.

Contrast this with the ordinary $L^2$ norm, $\norm{\theta}_{L^2}^2 =
\sum_k|\hat\theta(k)|^2$, which weights every Fourier mode equally and is
\emph{exactly conserved} by \eqref{eq:pde} (a direct consequence of $u$
being divergence-free: $\frac{d}{dt}\int\theta^2\dx = -2\int\theta\,
(u\cdot\nabla\theta)\dx = -\int u\cdot\nabla(\theta^2)\dx = \int
(\nabla\cdot u)\theta^2\dx = 0$). The mix norm instead weights the $k$-th
mode by $1/|k|^2$, penalising low wavenumbers (large-scale structure) and
suppressing high wavenumbers (fine-scale structure). A scalar field whose
$L^2$ energy sits at low $k$ (a smooth blob) therefore has \emph{large}
$\norm{\theta}_{H^{-1}}$, while the same $L^2$ energy redistributed to high
$k$ (fine filaments, the same substance stretched thin) has \emph{small}
$\norm{\theta}_{H^{-1}}$ -- even though $\norm{\theta}_{L^2}$ is unchanged.
This is the precise mathematical sense in which the mix norm captures
``mixedness'': a rearrangement that conserves $L^2$ mass while pushing
energy to finer scales is exactly what \eqref{eq:mixnorm} is designed to
detect.

\subsection{Enstrophy and the Stirring Budget}
\label{sec:enstrophy}

The ``stirring budget'' available to the velocity field is expressed as a
constraint on its \emph{enstrophy},
\begin{equation}
  \label{eq:enstrophy}
  \mathcal{E}(u) := \norm{\nabla u}_{L^2}^2
  = \int_\Torus |\nabla u_x|^2 + |\nabla u_y|^2 \dx.
\end{equation}
For a divergence-free planar field this coincides (up to the definition's
normalisation) with the $L^2$ norm of the scalar vorticity $\omega =
\partial_x u_y - \partial_y u_x$: writing $u$ in terms of a streamfunction
$\psi$ via $u = \nabla^\perp\psi = (-\partial_y\psi,\partial_x\psi)$ gives
$\omega = \Lap\psi$ and $\norm{\nabla u}_{L^2}^2 = \norm{\omega}_{L^2}^2$
after an integration by parts. Bounding $\norm{\nabla u}_{L^2}$ therefore
bounds the total squared vorticity of the flow -- physically, a cap on how
vigorously and how finely the flow is allowed to swirl. This is a stronger
constraint than capping the kinetic energy $\norm{u}_{L^2}$ alone: for fixed
enstrophy, a velocity field concentrated at high wavenumber (fine swirls)
uses up its budget faster than one concentrated at low wavenumber, since
$\norm{\nabla u}_{L^2}^2 = \sum_k 4\pi^2|k|^2|\hat u(k)|^2$.

\subsection{The Lin--Thiffeault--Doering Optimal Velocity}
\label{sec:ltd}

Among all incompressible velocity fields satisfying the enstrophy constraint
$\norm{\nabla u}_{L^2} \le F$, the one that minimises
$\frac{d}{dt}\norm{\theta}_{H^{-1}}^2$ at each instant is the
\emph{Lin--Thiffeault--Doering (LTD) optimal velocity}
\cite{LTD2011}:

\begin{equation}
  \label{eq:ltd}
  v = -\lapinv \Proj\!\left(\theta\,\nabla\lapinv\theta\right),
  \qquad
  u = F\,\frac{v}{\norm{\nabla v}_{L^2}}.
\end{equation}

Here $\Proj$ is the \emph{Leray projection} onto divergence-free vector
fields:

\begin{equation}
  \label{eq:leray}
  \Proj g = g - \nabla\lapinv(\nabla\cdot g).
\end{equation}

The Leray projection removes the curl-free (gradient) component of a vector
field $g$, leaving only its divergence-free component; this is precisely
what is required, since any admissible $u$ must itself be divergence-free.
The final rescaling by $\norm{\nabla v}_{L^2}^{-1}$ enforces the enstrophy
constraint exactly (with equality) at each instant -- intuitively, there is
never a reason to use less than the full stirring budget, since using more
of a fixed direction can only decrease $\norm{\theta}_{H^{-1}}^2$ faster.

\subsection{Derivation of the LTD Velocity}
\label{sec:derivation}

Formula~\eqref{eq:ltd} is not an ansatz; it is the exact solution of a
constrained optimisation problem, and deriving it makes transparent both
why the five-step algorithm of Section~\ref{sec:rhs} has the structure it
does, and when it can fail (Section~\ref{sec:saddle}).

Write $\phi := \lapinv\theta$ (so $\Lap\phi=\theta$), and recall from
Section~\ref{sec:mixnorm-def} that $\norm{\theta}_{H^{-1}}^2 =
-\int\theta\phi\dx$. Differentiating in time and using the symmetry of
$\Lap^{-1}$ (so that $\int\theta\dot\phi\dx = \int\dot\theta\phi\dx$),

\begin{equation}
  \frac{d}{dt}\norm{\theta}_{H^{-1}}^2
  = -2\int \dot\theta\,\phi \dx
  = 2\int (u\cdot\nabla\theta)\,\phi\dx,
\end{equation}

using \eqref{eq:pde}. Since $u$ is divergence-free, $\nabla\cdot(u\theta\phi)
= (u\cdot\nabla\theta)\phi + \theta(u\cdot\nabla\phi)$, and the divergence
integrates to zero on the torus, so $\int(u\cdot\nabla\theta)\phi\dx =
-\int\theta(u\cdot\nabla\phi)\dx$. Hence

\begin{equation}
  \label{eq:derivmixnorm}
  \frac{d}{dt}\norm{\theta}_{H^{-1}}^2
  = -2\int u\cdot\underbrace{(\theta\nabla\phi)}_{=:\,g}\dx
  = -2\int u\cdot g \dx.
\end{equation}

This is exactly the quantity $g=\theta\nabla\lapinv\theta$ computed in
Step~1 of Section~\ref{sec:rhs}. Because $u$ is divergence-free, the
curl-free part of $g$ contributes nothing to the pairing ($\int
u\cdot\nabla f \dx = -\int(\nabla\cdot u)f\dx = 0$ for any scalar $f$), so
$\int u\cdot g\dx = \int u\cdot\Proj g\dx$, and \eqref{eq:derivmixnorm}
becomes
\begin{equation}
  \label{eq:derivmixnorm2}
  \frac{d}{dt}\norm{\theta}_{H^{-1}}^2 = -2\int u\cdot\Proj g\dx.
\end{equation}

Now introduce $v := -\lapinv\Proj g$ as in \eqref{eq:ltd}, so that
$\Proj g = -\Lap v$, and
\begin{equation}
  \int u\cdot\Proj g\dx = -\int u\cdot\Lap v\dx
  = \int \nabla u : \nabla v \dx,
\end{equation}
again by integrating by parts (component-wise, using that both $u$ and $v$
are periodic and divergence-free). Substituting into
\eqref{eq:derivmixnorm2},
\begin{equation}
  \label{eq:derivmixnorm3}
  \frac{d}{dt}\norm{\theta}_{H^{-1}}^2 = -2\int \nabla u : \nabla v \dx.
\end{equation}

\begin{proposition}[Optimality of the LTD velocity]
\label{prop:ltd}
Among all divergence-free $u$ with $\norm{\nabla u}_{L^2}\le F$, the
quantity \eqref{eq:derivmixnorm3} is minimised by
$u = F\,v/\norm{\nabla v}_{L^2}$.
\end{proposition}
\begin{proof}
By the Cauchy--Schwarz inequality applied to the $L^2$ inner product of
gradients, $\int\nabla u:\nabla v\dx \le \norm{\nabla u}_{L^2}
\norm{\nabla v}_{L^2} \le F\norm{\nabla v}_{L^2}$, with equality iff $\nabla
u$ is a non-negative scalar multiple of $\nabla v$ with
$\norm{\nabla u}_{L^2}=F$, i.e.\ $u = F\,v/\norm{\nabla v}_{L^2}$ (up to an
additive constant, which vanishes on the mean-zero torus). This maximises
$\int\nabla u:\nabla v\dx$, hence minimises
$-2\int\nabla u:\nabla v\dx = \frac{d}{dt}\norm{\theta}_{H^{-1}}^2$.
\end{proof}

At this optimum, \eqref{eq:derivmixnorm3} evaluates to
$\frac{d}{dt}\norm{\theta}_{H^{-1}}^2 = -2F\norm{\nabla v}_{L^2} \le 0$,
confirming the monotonic decrease of the mix norm asserted informally in
Section~\ref{sec:mixnorm-def}, \emph{provided} $\norm{\nabla v}_{L^2}>0$.

\subsection{Degeneracy: the First-Order-Critical Case}
\label{sec:saddle}

The derivation above shows that the LTD velocity is well defined, and
strictly decreases the mix norm, exactly as long as $v\neq 0$, i.e.\ as
long as $\Proj g \neq 0$. If $\theta$ reaches a state where $g=\theta
\nabla\lapinv\theta$ is already curl-free (so that $\Proj g = 0$), then
\emph{no} divergence-free velocity, however large its enstrophy, can
decrease $\norm{\theta}_{H^{-1}}^2$ at that instant \cite{LTD2011}: such a
$\theta$ is a critical point of the mix-norm functional under
incompressible rearrangement. Lin, Thiffeault and Doering address this
degenerate case by instead seeking the incompressible velocity that
maximises the growth rate of the \emph{second} time derivative of the mix
norm, which reduces to an eigenvalue problem for the minimum eigenvalue of
a linear operator built from $\theta$ and $\varphi=\lapinv\theta$
\cite{LTD2011}. The numerical implementation here does not solve that
secondary eigenvalue problem; it instead monitors for proximity to the
degeneracy by comparing $\norm{\nabla v}_{L^2}$ against
$\norm{\lapinv g}_{L^2}$ and emitting a diagnostic warning if the former is
much smaller than the latter relative to a tolerance. In practice this
warning was not observed to trigger for the smooth, well-supported initial
data families used in Section~\ref{sec:idata}.

\begin{remark}
  The LTD velocity is the \emph{greedy} optimal: it minimises the
  instantaneous rate of change of the mix norm, not a global functional
  over time (Proposition~\ref{prop:ltd} is a pointwise-in-time statement).
  This instantaneous optimality does not by itself establish global
  optimality over a finite or infinite time horizon: Lin, Thiffeault and
  Doering themselves note that their local-in-time strategy leaves room
  for further improvement, and identify whether a global-in-time (optimal
  control) strategy could close the gap between their rigorous bounds and
  their simulations as a major open question \cite{LTD2011}.
  Section~\ref{sec:rate} below reports a numerical gap between the fitted
  decay-rate magnitude observed under the greedy LTD velocity and the
  support-size scaling suggested by the Iyer--Kiselev--Xu
  enstrophy-constrained lower bound \cite{Iyer2014}; Section~\ref{sec:resolution}
  shows that this gap is entangled with finite-resolution effects and
  should not, on its own, be read as proof of greedy sub-optimality.
\end{remark}

% ─────────────────────────────────────────────────────────────────────────────
\section{Numerical Methods}
\label{sec:numerics}

\subsection{Spectral Discretisation}

We discretise $[0,1]^2$ on an $N\times N$ grid with spacing $dx = 1/N$,
using the DFT convention

\begin{equation}
  \hat f[k_y,k_x]
  = \sum_{i=0}^{N-1}\sum_{j=0}^{N-1} f[i,j]\,e^{-2\pi i(ik_y+jk_x)/N}.
\end{equation}

Row index $i$ corresponds to the $y$-direction and column index $j$ to the
$x$-direction.
The centred wavenumbers are
$k_{\text{eff}}[j] = j - N\cdot\mathbf{1}[j>N/2]$,
giving the range $\{0,1,\ldots,N/2-1,-N/2+1,\ldots,-1\}$.

\paragraph{Derivative multipliers.}
Differentiation in physical space corresponds to multiplication in Fourier
space by
\begin{equation}
  \Dop{x}(k_x) = 2\pi i\,k_d(k_x),\qquad
  \Dop{y}(k_y) = 2\pi i\,k_d(k_y),
\end{equation}
where $k_d$ is identical to $k_{\text{eff}}$ except that the Nyquist mode
$k = N/2$ is set to zero (to avoid aliasing in first derivatives; see
Section~\ref{sec:aliasing}).

\paragraph{Inverse Laplacian.}
The discrete Fourier realisation of $\lapinv$ is the multiplier
\begin{equation}
  \lapinv(k_x,k_y)
  = \begin{cases}
      \dfrac{-1}{(2\pi)^2(k_x^2+k_y^2)} & (k_x,k_y)\neq(0,0),\\[6pt]
      0 & (k_x,k_y)=(0,0).
    \end{cases}
\end{equation}
The Nyquist mode is \emph{not} zeroed here, since the Laplacian is a
second-order (even) operator for which the sign ambiguity of the Nyquist
frequency does not arise.

\paragraph{Mix-norm weight.}
\begin{equation}
  \Linv(k_x,k_y)
  = \frac{1}{2\pi|k|}
  = \sqrt{-\lapinv(k_x,k_y)}.
\end{equation}

\paragraph{Parseval identity.}
For the unitary DFT normalisation used here,
$\norm{\hat f}_{\ell^2} = N\norm{f}_{\ell^2}$, so

\begin{equation}
  \norm{f}_{L^2} = \frac{\norm{\hat f}_{\ell^2}}{N^2}.
\end{equation}

This conversion factor is used throughout whenever a quantity computed in
Fourier space (via $\ell^2$ norms of coefficient arrays) must be compared
against, or combined with, a continuum $L^2$ quantity.

\subsection{Pseudo-Spectral Evaluation of the RHS}
\label{sec:rhs}

At each RK45 stage the following five steps compute $\partial_t\hat\theta$.
The alternation between Fourier space (for the linear operators $\nabla$
and $\lapinv$, applied as cheap pointwise multiplications) and physical
space (for the pointwise \emph{products} that make up the nonlinearity)
is the defining feature of the pseudo-spectral method: representing a
product as a convolution directly in Fourier space would cost $O(N^4)$
per evaluation, whereas transforming to physical space, multiplying
pointwise, and transforming back costs only $O(N^2\log N)$ via the FFT.

\begin{enumerate}
  \item \textbf{Compute} $g = \theta\,\nabla\lapinv\theta$ in physical space
        (pseudo-spectral):
        \[
          \hat g_x = \mathcal{F}\!\left[\theta\cdot\mathcal{F}^{-1}
            (\Dop{x}\cdot\lapinv\cdot\hat\theta)\right],
          \quad
          \hat g_y = \mathcal{F}\!\left[\theta\cdot\mathcal{F}^{-1}
            (\Dop{y}\cdot\lapinv\cdot\hat\theta)\right].
        \]
        This is exactly the quantity $g$ appearing in the derivation of
        Section~\ref{sec:derivation}.

  \item \textbf{Project} onto divergence-free fields (Leray):
        \[
          \widehat{\Proj g}_x = \hat g_x
            - \Dop{x}\cdot\lapinv
              \cdot(\Dop{x}\cdot\hat g_x+\Dop{y}\cdot\hat g_y),
        \]
        and similarly for the $y$-component.

  \item \textbf{Compute} $v = -\lapinv\Proj g$:
        \[
          \hat v_x = -\lapinv\cdot\widehat{\Proj g}_x,\qquad
          \hat v_y = -\lapinv\cdot\widehat{\Proj g}_y.
        \]

  \item \textbf{Normalise} to enforce the enstrophy constraint via Parseval:
        \[
          \norm{\nabla v}_{L^2}
          = \frac{1}{N^2}\sqrt{
              \norm{\Dop{x}\hat v_x}_{\ell^2}^2
              +\norm{\Dop{x}\hat v_y}_{\ell^2}^2
              +\norm{\Dop{y}\hat v_x}_{\ell^2}^2
              +\norm{\Dop{y}\hat v_y}_{\ell^2}^2},
          \quad
          u = F\,\frac{v}{\norm{\nabla v}_{L^2}}.
        \]
        This is Proposition~\ref{prop:ltd} applied numerically: $v$ is
        rescaled to have exactly enstrophy $F$.

  \item \textbf{Advect}:
        \[
          \frac{d\hat\theta}{dt}
          = -\mathcal{F}\!\left[
              u_x\cdot\mathcal{F}^{-1}(\Dop{x}\cdot\hat\theta)
              + u_y\cdot\mathcal{F}^{-1}(\Dop{y}\cdot\hat\theta)
            \right].
        \]
\end{enumerate}

Each RHS evaluation therefore requires ten $N\times N$ FFTs (three forward
and seven inverse transforms, counting the intermediate transforms in
Steps 1--5), each
costing $O(N^2\log N)$, for a total cost of $O(N^2\log N)$ per stage --
independent of $N$ only up to the logarithmic factor, which is why doubling
$N$ from $32$ to $64$ increases per-step cost by roughly a factor of
$4$--$5$ (four times the grid points, times a modestly larger logarithm),
consistent with the timings reported in Section~\ref{sec:resolution}.

\subsection{Aliasing and the Absence of Explicit Dealiasing}
\label{sec:aliasing}

The nonlinearity in Step~1 above is a \emph{pointwise product} of two
fields, each itself resolved on $N$ Fourier modes per direction. In exact
arithmetic such a product generates Fourier content up to wavenumber $N-1$
(twice the input bandwidth), but the discrete grid can only represent
wavenumbers up to $N/2$; the un-representable high-wavenumber content is
\emph{aliased} back onto low wavenumbers, contaminating the low modes with
spurious energy. The classical remedy in pseudo-spectral codes is the
``$2/3$-rule'': pad the field with zeros to $3N/2$ modes before
transforming, multiply, and truncate back to the original $N/2$
bandwidth, exactly eliminating the aliased content at the cost of roughly
doubling the FFT size.

The implementation does \emph{not} apply the $2/3$-rule (nor any other
explicit dealiasing filter beyond zeroing the Nyquist mode for first
derivatives, which addresses a different, sign-ambiguity issue, not
aliasing). Instead, the resolution-check event of
Section~\ref{sec:rescheck} provides an \emph{a posteriori} diagnostic of
loss of resolution: because under-resolution -- aliasing error among other
possible causes -- corrupts the conserved $L^p$ norms of $\theta$, its
cumulative effect is what the resolution check is designed to detect, and
integration is halted the moment that effect becomes non-negligible. This
monitor does not \emph{prevent} aliasing error from occurring and is not a
substitute for explicit dealiasing; it only bounds how long the simulation
is allowed to run once accumulated resolution loss becomes detectable. The
consequence -- a stopping time that
depends on resolution, and hence a fitted mixing-rate exponent that depends
on resolution -- is examined directly in Section~\ref{sec:resolution} and is
relevant to the deviation from the Iyer--Kiselev--Xu benchmark $a^{-1}$
scaling discussed in Section~\ref{sec:rate}.

\subsection{Time Integration}
\label{sec:timeint}

The ODE is integrated with an explicit, adaptive-step embedded
Runge--Kutta pair of Dormand--Prince type (order 5(4)), with relative and
absolute error tolerances of $10^{-6}$ and $10^{-8}$ respectively: the
solution is advanced with a 5th-order formula while a simultaneously
computed 4th-order formula provides a local error estimate, and the step
size is adjusted so that this estimate stays within tolerance. No
dealiasing-related CFL restriction is imposed explicitly; the adaptive step
controller implicitly shrinks the step size as $\theta$ develops sharper
gradients under stirring, which is the mechanism by which the integrator
remains stable as filamentation proceeds, right up until the resolution
check intervenes.

Because the solver requires a real-valued state vector, the complex
Fourier array $\hat\theta\in\mathbb{C}^{N^2}$ is stored as a real vector
$y = [\operatorname{Re}(\hat\theta_{\text{flat}}),\,
       \operatorname{Im}(\hat\theta_{\text{flat}})]\in\mathbb{R}^{2N^2}$.

\subsection{Resolution Check and Stopping Criterion}
\label{sec:rescheck}

Under incompressible advection all $L^p$ norms of $\theta$ are conserved.
We monitor the event function

\begin{equation}
  e(t) = \max\!\left(
    |\norm{\theta}_{L^2}-1|,\;
    \left|\frac{\norm{\theta}_{L^4}}{\norm{\theta_0}_{L^4}}-1\right|,\;
    \left|\frac{\norm{\theta}_{L^8}}{\norm{\theta_0}_{L^8}}-1\right|
  \right) - \varepsilon,
  \quad \varepsilon = 10^{-3}.
\end{equation}

Integration stops when $e(t)$ crosses zero upward, indicating that the
grid can no longer resolve the fine-scale features of $\theta$. Using three
different exponents ($p=2,4,8$) rather than a single one is deliberate: the
higher-order norms are more sensitive to sharp, localised features (large
$|\theta|$ concentrated on a small area) than $L^2$, so drift is typically
detected first in $L^8$, giving an early warning before the coarser $L^2$
budget is visibly affected -- consistent with the observation in
Section~\ref{sec:conservation} that all three norms begin drifting at
essentially the same instant, since by that point the under-resolved
filaments are thin enough to affect every norm simultaneously.

% ─────────────────────────────────────────────────────────────────────────────
\section{Initial Data}
\label{sec:idata}

Four families of initial data are used, all $L^2$-normalised
($\norm{\theta_0}_{L^2}=1$) and parametrised by a scale parameter
$a\in(0,1)$ that controls the support size.
After construction the support is circularly shifted to the centre of $[0,1]^2$
so that the mixing dynamics are not biased by proximity to the periodic
boundary.

\begin{center}
\begin{tabular}{lll}
  \toprule
  Family & Support & Description \\
  \midrule
  I: Sine bump & $[0,a]^2$ & $\sin(2\pi x/a)\sin(2\pi y/a)$ \\
  II: Diagonal pair & Two sub-squares of $[0,a]^2$ & Antisymmetric pair,
                                                       broken symmetry \\
  III: Strip & $[0,a]\times[0,a/2]$ & Aspect-ratio $2{:}1$ strip \\
  IV: Trigonometric polynomial & $[0,1]^2$ &
        $\sum_{k=1}^{2}c_k\sin(2k\pi x/a)\sin(2k\pi y/a)$ \\
  \bottomrule
\end{tabular}
\end{center}

All results in Section~\ref{sec:results} use the sine-bump family (Family
I), the simplest of the four, so that the effect of $a$ on the mixing rate
is not confounded by the initial data's internal geometry; the other three
families are used for testing sensitivity to the initial condition's shape
(a direction identified as future work in Section~\ref{sec:conc}).

Figure~\ref{fig:idata} shows all four functions for $a=0.5$.

\begin{figure}[H]
  \centering
  \includegraphics[width=\linewidth]{figures/fig1_initial_conditions.pdf}
  \caption{%
    The four families of initial data $\theta_0$ with scale parameter $a=0.5$
    on the $N=64$ grid.
    Colour encodes the scalar value (red $>0$, blue $<0$, white $= 0$).
  }
  \label{fig:idata}
\end{figure}

% ─────────────────────────────────────────────────────────────────────────────
\section{Simulation Results}
\label{sec:results}

All results in this section use $N=64$ spectral modes per direction and
enstrophy constraint $F=1$. Each simulation starts at $t=0$ and runs until
the resolution check event fires. Every figure and number in this section
was produced by the accompanying code and independently re-verified: a
clean rerun of the pipeline reproduces the values originally recorded in
Table~\ref{tab:rates} to four decimal places in every entry, confirming
that the pipeline is deterministic and reproducible end to end for the
tested software environment.

\subsection{Spatial Evolution of the Scalar Field}
\label{sec:snapshots}

Figure~\ref{fig:snaps} shows six snapshots of the scalar field $\theta(x,y,t)$
for $a=0.5$, sampled evenly between $t=0$ and the stopping time $t=1.30$.
The LTD velocity progressively stirs $\theta$ into finer and finer filaments,
consistent with exponential decay of the mix norm. Qualitatively, the
sequence illustrates the classical picture of chaotic advection: smooth
initial structure is repeatedly stretched and folded, and the length of the
scalar's interface grows while its cross-sectional width shrinks, which is
precisely the geometric process that pushes $L^2$ energy to higher
wavenumbers and hence drives $\norm{\theta}_{H^{-1}}$ down.

\begin{figure}[H]
  \centering
  \includegraphics[width=\linewidth]{figures/fig2_snapshots.pdf}
  \caption{%
    Snapshots of the passive scalar $\theta(x,y,t)$ at six equally spaced
    times for $a=0.5$, $N=64$, $F=1$.
    The LTD velocity continuously stretches and folds the scalar into
    progressively finer structures.
  }
  \label{fig:snaps}
\end{figure}

\subsection{Mix Norm Decay}
\label{sec:mixnorm}

Figure~\ref{fig:mixnorm} shows $\log(\norm{\theta(t)}_{H^{-1}} /
\norm{\theta_0}_{H^{-1}})$ for $a=0.5$.
After a brief initial transient the decay is well approximated by a linear
fit (dashed), confirming approximately exponential decay:

\begin{equation}
  \norm{\theta(t)}_{H^{-1}} \approx \norm{\theta_0}_{H^{-1}}\,e^{-rt},
  \qquad r>0.
\end{equation}

Here $r$ denotes the positive decay-rate magnitude, so that the fitted
slope of $\log\norm{\theta}_{H^{-1}}$ versus $t$ equals $-r$ and the mixing
timescale is $\tau=1/r$. Fitting the slope over the \emph{entire}
trajectory gives slope $\approx -0.394$ (decay-rate magnitude $r \approx
0.394$); fitting only the \emph{last two-thirds} (the methodology used for
Table~\ref{tab:rates} and Figure~\ref{fig:rate}, which discards the initial
transient before the exponential regime is established) gives the steeper
slope $\approx -0.440$ ($r \approx 0.440$) quoted in Table~\ref{tab:rates}.
The two numbers are not in conflict -- they measure the average slope over
different portions of a curve that is only \emph{approximately} log-linear
-- but the sensitivity of the fitted rate to the choice of fitting window is
itself informative: it shows that the decay, while broadly exponential, has
curvature that a single global rate does not fully capture, most visible
during the initial transient before the flow has had time to establish
fully developed filamentation.

\begin{figure}[H]
  \centering
  \includegraphics[width=0.75\linewidth]{figures/fig3_mixnorm_single.pdf}
  \caption{%
    Log of the $H^{-1}$ mix norm vs.\ time for $a=0.5$.
    The dashed line is the least-squares linear fit over the full trajectory,
    giving slope $-0.3944$.
  }
  \label{fig:mixnorm}
\end{figure}

\subsection{Conservation Law Verification}
\label{sec:conservation}

Figure~\ref{fig:lp} confirms that the $L^2$, $L^4$ and $L^8$ norms remain
$\approx 1$ throughout the integration.
The norms begin to drift near $t=1.3$, at which point the resolution check
terminates the simulation.
The maximum $L^2$ drift observed is $\max_t|\norm{\theta}_{L^2}-1|
\approx 3.15\times10^{-5}$, two orders of magnitude within the tolerance
$\varepsilon=10^{-3}$ -- indicating that the resolution check fires
promptly once drift becomes appreciable, rather than allowing substantial
error to accumulate silently beforehand.

\begin{figure}[H]
  \centering
  \includegraphics[width=0.72\linewidth]{figures/fig4_lp_norms.pdf}
  \caption{%
    $L^p$ norms normalised to their initial values for $a=0.5$.
    All three norms should stay $\approx 1$; a sustained drift triggers the
    resolution-check event that stops the integration.
  }
  \label{fig:lp}
\end{figure}

\subsection{Dependence on Scale Parameter $a$}
\label{sec:sweep}

We run the simulation for eight values of $a$ in the range
$[0.5, 0.9375]$ in steps of $1/16$ (matching the range used in the original
study).
Figure~\ref{fig:sweep} shows the log mix norm for all eight runs.
Smaller $a$ (narrower initial support) leads to a faster initial decay
but also runs into resolution limitations earlier, since a smaller initial
blob is stretched into sub-grid-scale filaments sooner.

\begin{figure}[H]
  \centering
  \includegraphics[width=0.85\linewidth]{figures/fig5_mixnorm_sweep.pdf}
  \caption{%
    $\log(\norm{\theta}_{H^{-1}}/\norm{\theta_0}_{H^{-1}})$ vs.\ time
    for $a\in\{0.5, 0.5625,\ldots,0.9375\}$ (colour coded from dark to light).
    Each run terminates when the resolution check fires.
  }
  \label{fig:sweep}
\end{figure}

\subsection{Mixing Timescale vs.\ Scale Parameter}
\label{sec:rate}

To quantify the mixing rate for each $a$, we fit a line to
$\log\norm{\theta}_{H^{-1}}$ over the \emph{last two-thirds} of each
trajectory (discarding the initial transient); the fitted slope is
$-r(a)$ with decay-rate magnitude $r(a)>0$, and the mixing timescale is
$\tau(a) = 1/r(a)$.

Table~\ref{tab:rates} summarises the results.
Figure~\ref{fig:rate} shows $\tau(a)$ plotted against $a$ together with a
linear fit (dashed red).

\begin{table}[H]
  \centering
  \caption{%
    Fitted decay-rate magnitudes $r$ and mixing timescales $\tau=1/r$
    for each scale parameter $a$.
    ($N=64$, $F=1$, initial data: sine-bump family.)
    These values were regenerated independently for this report and match
    the values originally recorded for this pipeline to four decimal
    places in every entry.
  }
  \label{tab:rates}
  \begin{tabular}{ccccc}
    \toprule
    $a$ & $t_{\text{stop}}$ & $\norm{\theta(t_{\text{stop}})}_{H^{-1}}$
        & $r$ & $\tau = 1/r$ \\
    \midrule
    0.5000 & 1.30 & 0.6058 & $0.4400$ & 2.27 \\
    0.5625 & 1.60 & 0.5831 & $0.3858$ & 2.59 \\
    0.6250 & 1.95 & 0.5658 & $0.3342$ & 2.99 \\
    0.6875 & 2.30 & 0.5627 & $0.2863$ & 3.49 \\
    0.7500 & 2.70 & 0.5642 & $0.2436$ & 4.11 \\
    0.8125 & 3.00 & 0.5862 & $0.2074$ & 4.82 \\
    0.8750 & 3.25 & 0.6220 & $0.1741$ & 5.74 \\
    0.9375 & 3.20 & 0.6920 & $0.1423$ & 7.03 \\
    \bottomrule
  \end{tabular}
\end{table}

Two features of Table~\ref{tab:rates} deserve comment beyond the monotone
growth of $\tau(a)$. First, $t_{\text{stop}}$ grows with $a$ -- larger
initial blobs survive longer before the resolution check fires, simply
because they take longer to be stretched down to the grid scale. Second,
$\norm{\theta(t_{\text{stop}})}_{H^{-1}}$ is \emph{not} monotone in $a$: it
decreases from $a=0.5$ to a minimum around $a\approx0.6875$--$0.75$, then
increases again toward $a=0.9375$. This non-monotonicity reflects a
competition between two effects that both depend on $a$: smaller blobs mix
faster per unit time (favouring small final mix norm) but are also stopped
sooner by the resolution check (favouring a less-mixed final state); the
minimum marks where these two effects roughly balance.

\begin{figure}[H]
  \centering
  \includegraphics[width=0.65\linewidth]{figures/fig6_mixing_rate.pdf}
  \caption{%
    Mixing timescale $\tau(a) = 1/r(a)$ vs.\ $a$.
    The fitted decay-rate magnitude scales as $r \sim a^{-1.78}$ (power-law
    fit to $\log r$ vs.\ $\log a$), compared with the $a^{-1}$ scaling
    characteristic of the Iyer--Kiselev--Xu enstrophy-constrained lower
    bound~\cite{Iyer2014}.
    The dashed red line is a linear fit to $\tau$ vs.\ $a$.
  }
  \label{fig:rate}
\end{figure}

\begin{remark}
  The observed scaling $r\sim a^{-1.78}$ deviates from the $a^{-1}$
  dependence characteristic of the rigorous exponential lower bound of
  Iyer, Kiselev and Xu for enstrophy-constrained mixing of compactly
  supported two-dimensional data of support area $O(a^2)$
  \cite{Iyer2014}. That result is a lower bound on how fast the
  $H^{-1}$ norm can possibly decay, not an exact prediction for the
  dynamically chosen LTD decay rate, so a discrepancy between the two is
  not automatically evidence of any single mechanism (finite resolution,
  sub-optimality of the greedy strategy, or both). Longer integration
  times (larger $N$ to delay the resolution event) may be needed to
  isolate the asymptotic regime. Section~\ref{sec:resolution} examines
  this by comparing the fitted exponent at two different resolutions.
\end{remark}

% ─────────────────────────────────────────────────────────────────────────────
\section{Numerical Verification and Resolution Sensitivity}
\label{sec:resolution}

This section is not present in Gautam Iyer's earlier MATLAB implementation;
it was added to (i) independently verify the reproducibility of
Table~\ref{tab:rates}, and (ii) directly test the Remark of
Section~\ref{sec:rate} -- that the deviation from the $a^{-1}$ scaling may
be, in part, a finite-resolution artefact -- by rerunning the same
eight-value sweep of Section~\ref{sec:sweep} at a coarser resolution,
$N=32$, and comparing the fitted results.

\subsection{Reproducibility at $N=64$}

Regenerating Table~\ref{tab:rates} from a clean run of the simulation
reproduces every entry to four decimal places: $t_{\text{stop}}$,
$\norm{\theta(t_{\text{stop}})}_{H^{-1}}$, and $r$ all match exactly,
and the fitted power law recovers $r\sim a^{-1.777}$, matching the
$a^{-1.78}$ quoted in Section~\ref{sec:rate} to the precision reported
there. This confirms deterministic reproducibility of the RK45
integration, the resolution-check event, and the least-squares fitting
procedure for the tested software environment and parameter configuration
used in this study; it does not establish independence from library
versions or execution environments, which were not varied here.

\subsection{Resolution Dependence of the Fitted Exponent}

Table~\ref{tab:resolution} compares the $N=64$ results of
Table~\ref{tab:rates} against the same eight-value $a$-sweep run at
$N=32$, using identical initial data, enstrophy constraint, and fitting
procedure.

\begin{table}[H]
  \centering
  \caption{%
    Resolution sensitivity: mixing timescale $\tau(a)=1/r(a)$ at
    $N=32$ versus $N=64$, same $a$-sweep and fitting procedure.
  }
  \label{tab:resolution}
  \begin{tabular}{ccccc}
    \toprule
    $a$ & $t_{\text{stop}}$ ($N{=}32$) & $\tau$ ($N{=}32$)
        & $t_{\text{stop}}$ ($N{=}64$) & $\tau$ ($N{=}64$) \\
    \midrule
    0.5000 & 0.65 & 2.83 & 1.30 & 2.27 \\
    0.5625 & 0.75 & 3.23 & 1.60 & 2.59 \\
    0.6250 & 1.00 & 3.49 & 1.95 & 2.99 \\
    0.6875 & 1.20 & 3.97 & 2.30 & 3.49 \\
    0.7500 & 1.45 & 4.51 & 2.70 & 4.11 \\
    0.8125 & 1.75 & 5.23 & 3.00 & 4.82 \\
    0.8750 & 1.90 & 6.44 & 3.25 & 5.74 \\
    0.9375 & 2.00 & 8.29 & 3.20 & 7.03 \\
    \bottomrule
  \end{tabular}
\end{table}

Two observations follow directly from Table~\ref{tab:resolution}. First,
$t_{\text{stop}}$ at $N=32$ is uniformly and substantially smaller than at
$N=64$ -- roughly half, across every value of $a$ -- confirming that coarser
grids hit the resolution limit sooner, exactly as expected since a coarser
grid can resolve fewer decades of filamentation before under-resolution
contaminates the conserved $L^p$ norms. Second, the fitted power-law
exponent itself changes with resolution: $r\sim a^{-1.61}$ at $N=32$
versus $r\sim a^{-1.78}$ at $N=64$. Since both exponents are fitted
from trajectories that are truncated early by the same resolution-check
mechanism, and the truncation point itself moves with $N$
(Section~\ref{sec:aliasing}), this change from approximately $-1.61$ to
$-1.78$ demonstrates that the fitted exponent is resolution-dependent and
not yet numerically converged. Because the higher-resolution estimate
moves \emph{farther} from the benchmark $-1$ scaling rather than closer to
it, these two resolutions alone do not establish that finite-resolution
effects explain the discrepancy from the Iyer--Kiselev--Xu benchmark
scaling \cite{Iyer2014}: they show only that at least two resolutions are
insufficient to determine the limiting exponent, and cannot distinguish
whether higher resolution would approach $-1$, remain near $-1.8$, or
behave otherwise. Disentangling genuine greedy-strategy behaviour from
finite-resolution truncation would require runs at substantially higher
$N$ (and correspondingly finer, possibly implicit or exponential,
time-stepping to keep the cost tractable), which is identified as future
work in Section~\ref{sec:conc}.

% ─────────────────────────────────────────────────────────────────────────────
\section{Conclusions}
\label{sec:conc}

We have numerically reproduced and extended Gautam Iyer's earlier MATLAB
implementation of optimal passive-scalar mixing, adding a first-principles
derivation of the LTD optimal velocity (Section~\ref{sec:derivation}), an
explicit account of the numerical scheme's treatment (and non-treatment)
of aliasing error (Section~\ref{sec:aliasing}), and a resolution-sensitivity
study not present in that earlier implementation (Section~\ref{sec:resolution}).
The Python pipeline is deterministically reproducible for the tested
software environment: a clean rerun reproduces the values reported here to
four decimal places in every entry.

The main numerical findings are:
\begin{itemize}
  \item The $H^{-1}$ mix norm decays approximately exponentially under the
        LTD optimal velocity, consistent with the theoretical predictions of
        \cite{LTD2011}, and consistent with the monotonic-decrease property
        established rigorously in Proposition~\ref{prop:ltd}.
  \item The mixing timescale $\tau(a) = 1/r(a)$ grows with $a$,
        confirming that scalars with wider initial support mix more slowly,
        while the mix norm reached at the stopping time is non-monotone in
        $a$, reflecting a trade-off between faster mixing and earlier
        resolution loss for smaller $a$.
  \item The observed power-law scaling $r\sim a^{-1.78}$ deviates from the
        $a^{-1}$ dependence characteristic of the Iyer--Kiselev--Xu
        enstrophy-constrained lower bound; the resolution-sensitivity
        study of Section~\ref{sec:resolution} shows this exponent is
        itself resolution-dependent and not yet numerically converged
        ($a^{-1.61}$ at $N=32$ vs.\ $a^{-1.78}$ at $N=64$, moving
        \emph{away} from $-1$ rather than toward it), so these two
        resolutions cannot determine how much of the deviation, if any,
        reflects a genuine feature of the greedy strategy versus a
        finite-resolution artefact.
  \item The $L^p$ conservation laws provide a robust, physically motivated
        \emph{a posteriori} diagnostic of resolution loss and a stopping
        criterion (Section~\ref{sec:aliasing}), but they do not prevent
        aliasing and are not a substitute for explicit dealiasing: the
        simulation instead halts automatically once under-resolution
        becomes detectable, rather than allowing it to accumulate
        silently.
\end{itemize}

Future work could explore:
\begin{itemize}
  \item Higher resolution ($N=128$ or $N=256$), ideally combined with
        explicit $2/3$-rule dealiasing, to separate genuine mixing-rate
        sub-optimality from finite-resolution truncation effects, and to
        access the true asymptotic regime before the resolution check
        intervenes.
  \item Alternative initial data (e.g.\ random fields, or the
        diagonal-pair, strip, and trigonometric-polynomial families of
        Section~\ref{sec:idata} not swept in Section~\ref{sec:results})
        and their effect on the mixing rate.
  \item Comparison with non-greedy optimal strategies, e.g.\ velocities
        that account for the effect of stirring over a finite time horizon
        rather than instantaneously, to directly probe the greedy-vs-global
        optimality question raised in \cite{LTD2011}.
  \item A systematic study of the first-order degeneracy identified in
        Section~\ref{sec:saddle}, including initial data explicitly
        engineered to approach $\Proj g \approx 0$.
\end{itemize}

% ─────────────────────────────────────────────────────────────────────────────
\bibliographystyle{plainnat}
\begin{thebibliography}{9}

\bibitem{LTD2011}
Z.~Lin, J.-L.~Thiffeault, and C.~R.~Doering,
\newblock Optimal stirring strategies for passive scalar mixing,
\newblock \emph{Journal of Fluid Mechanics}, 675:465--476, 2011.

\bibitem{Iyer2014}
G.~Iyer, A.~Kiselev, and X.~Xu,
\newblock Lower bounds on the mix norm of passive scalars advected by
  incompressible enstrophy-constrained flows,
\newblock \emph{Nonlinearity}, 27(5):973--985, 2014.

\bibitem{Trefethen2000}
L.~N.~Trefethen,
\newblock \emph{Spectral Methods in MATLAB},
\newblock SIAM, 2000.

\bibitem{Boyd2001}
J.~P.~Boyd,
\newblock \emph{Chebyshev and Fourier Spectral Methods},
\newblock 2nd edition, Dover Publications, 2001.

\end{thebibliography}

\end{document}
